\begin{lemma}
Let \(I\subseteq[n]\) satisfy
\[
|I|=k=\noderef{TwoBiteNaturalIndependenceScale}(1+\varepsilon,n),
\]
and assume \(0\le\varepsilon_1\le1\),
\(\noderef{ParameterHierarchyEventualComparisons}
(\varepsilon,\varepsilon_1,\varepsilon_2,n_0)\), \(n>n_0\), and
\[
m=\noderef{TwoBiteNaturalReducedVertexCount}(n),\qquad
p=\noderef{TwoBiteEdgeProbability}\!\left(\frac12,n\right).
\]
Suppose that every sample \(\omega\) satisfying both
\[
\noderef{TwoBiteRegularityEvent}(k,\varepsilon,\varepsilon_1,\varepsilon_2,
\omega)
\]
and
\[
\noderef{TwoBiteProjectionSizeEvent}(I,k,\ell_R,\ell_B,\omega)
\]
also satisfies
\[
\noderef{ProjectionOpenPairFunction}(\ell_R,\ell_B,k,p,n)
  -10\varepsilon_1k^2
\le
\left|\noderef{TwoBiteStagedOpenPairs}(\omega,\varepsilon,I)\right|.
\tag{1}
\]
Then
\[
\noderef{TwoBiteConditionalProbability}\!\left(n,m,p;\,
I\hbox{ is independent in }\noderef{TwoBiteFinalGraph}(\omega)_0
\text{ and }
\noderef{TwoBiteRegularityEvent}(k,\varepsilon,\varepsilon_1,\varepsilon_2,
\omega)
\;\middle|\;
\noderef{TwoBiteProjectionSizeEvent}(I,k,\ell_R,\ell_B,\omega)
\right)
\le
\exp\!\left(
-p f(\ell_R,\ell_B)
+8\sqrt{\varepsilon_1}\,pk^2
+10\varepsilon_1pk^2
+4\log k\right),
\]
where \(f=\noderef{ProjectionOpenPairFunction}(\ell_R,\ell_B,k,p,n)\).
\end{lemma}
\begin{proof}
Let \(S=\mathcal S_{I,\ell_R,\ell_B}\).  If \(\Pr(S)=0\), the definition of
\(\noderef{TwoBiteConditionalProbability}\) makes the conditional probability
equal to \(0\), so the asserted upper bound is immediate for the nonnegative
bounds used here.  Assume \(\Pr(S)>0\).

The hierarchy hypotheses and the identity
\(p=\noderef{TwoBiteEdgeProbability}(\frac12,n)
=\frac12\sqrt{\log n/n}\) put \(p\) in the strict Bernoulli range
\(0<p<1\): the eventual range has \(\log n>0\), so the displayed formula
gives \(p>0\), while
\(\noderef{ParameterHierarchyEventualComparisons}\) gives
\(p\le1/2<1\).
\(\noderef{FixedSetExposureCellProductLaw}\), applied to this projection-size
cell \(S\), supplies a finite partition of \(S\) into exposure histories
\(C_h\).  On each \(C_h\) there is a fixed terminal candidate set \(U_h\) and
a fixed order in which the coordinates of \(U_h\) are exposed.  It also
supplies the recorded set \(R_h\) with
\[
e\in R_h
\Longleftrightarrow
e\in\noderef{TwoBitePreTerminalRecordedPairs}(\omega,\varepsilon,I)
\qquad(\omega\in C_h),
\tag{H1}
\]
and the full terminal-universe identity
\[
e\in U_h
\Longleftrightarrow
e\in\noderef{TwoBiteTerminalCoordinateUniverse}(m)\hbox{ and }e\notin R_h.
\tag{H2}
\]
The coordinates of \(U_h\) have product Bernoulli-\(p\) law after
conditioning on \(C_h\).  The staged-open set is not fixed on the cell; for
each
\(\omega\in C_h\) it is an adaptive subset
\[
T(\omega)=\noderef{TwoBiteStagedOpenPairs}(\omega,\varepsilon,I)
        \subseteq U_h.
\]
The prefix-safety part of the helper says that, when a coordinate of \(U_h\)
is reached in the terminal order, whether it currently belongs to
\(T(\omega)\) is already determined from the history and earlier terminal
values.  Its own edge value is still a fresh Bernoulli-\(p\) factor.  The
helper also says that a bound proved uniformly for every history on the event
\[
\mathcal B_I\cap\mathcal R
\]
averages to the identical bound conditioned only on \(S\).  It is therefore
enough to fix one compatible history \(h\) and prove the displayed
exponential bound for
\(\Pr(\mathcal B_I\cap\mathcal R\mid C_h)\).

Let \(\mathcal B_I\) be the event that \(I\) is independent in
\(\noderef{TwoBiteFinalGraph}\).  If
\(\Pr(\mathcal B_I\cap\mathcal R\mid C_h)=0\), there is nothing to prove.
Otherwise some sample in \(C_h\) satisfies \(\mathcal B_I\cap\mathcal R\).
For any such sample \(\omega\), (1) applies because the history lies inside
\(S\), and gives the samplewise lower bound
\[
|T_R(\omega)|+|T_B(\omega)|=|T(\omega)|
\ge f(\ell_R,\ell_B)-10\varepsilon_1k^2. \tag{2}
\]
Here \(T_R(\omega)\) and \(T_B(\omega)\) are the red and blue parts of the
adaptive staged-open set \(T(\omega)\).

If \(\mathcal B_I\cap\mathcal R\) occurs, then every present base edge among
the pairs of \(T_R(\omega)\cup T_B(\omega)\) must be deleted by an
opposite-color witness.  A present red staged-open coordinate cannot survive
into the final graph, or two vertices of \(I\) would be adjacent.  The
same-color deletion alternative in which this coordinate is the latest
monochromatic edge is excluded by
\(\noderef{TwoBiteProjectionPairSameColorClosed}\), and the already-revealed
alternative is excluded by
\(\noderef{TwoBiteProjectionPairTouched}\).  Hence a present red coordinate
must have a blue witness; the color-swapped statement holds for present blue
coordinates.

Let \(P_R(\omega)\subseteq T_R(\omega)\) and
\(P_B(\omega)\subseteq T_B(\omega)\) be the present red and blue base edges
among the staged-open pairs, and write
\[
u_R=|P_R(\omega)|,\qquad u_B=|P_B(\omega)|.
\]
There are at most \(k^2\) possible values for \(u_R\) and at most \(k^2\)
possible values for \(u_B\), because every staged-open coordinate comes from
projected pairs of a \(k\)-element set.  The final summation over the two
cardinalities therefore contributes at most \(k^4\).

The exceptional choices must be fixed before the same-color Bernoulli
variables they constrain are exposed.  This is supplied by
\(\noderef{FixedSetAdaptiveExceptionalCandidates}\).  When \(u_R\ge u_B\),
we reveal the blue terminal candidates first in the fixed history order.  Fix
one possible blue realization and encode it by a completed sample
\(\omega^\ast\in C_h\); only the injection, the recorded table, and the blue
terminal values of \(\omega^\ast\) are used.  We instantiate
\(\noderef{FixedSetAdaptiveExceptionalCandidates}\) with the history cell
\(C_h\), the recorded set \(R_h\), and the terminal set \(U_h\).  The
stronger hypotheses it requires are exactly (H1), (H2), and the
staged-open containment \(T(\omega)\subseteq U_h\), all supplied by
\(\noderef{FixedSetExposureCellProductLaw}\).  The helper defines a red
candidate set
\[
E_R^\ast\subseteq U_h
\]
from this information alone: a red coordinate belongs to \(E_R^\ast\) exactly
when it has a large, medium, or small blue opposite-color witness in the
first-color realization.  In particular \(E_R^\ast\) is fixed before any red
terminal Bernoulli value is exposed.  For every regular completion with this
same blue realization,
\[
|E_R^\ast|\le 3\varepsilon_1k^2, \tag{3}
\]
and if \(\mathcal B_I\) occurs, every present red staged-open coordinate of
that completion lies in \(E_R^\ast\).  The containment is not an after-the-fact
choice of an exceptional set: it follows by applying
\(\noderef{FixedSetExceptionalWitnessClassification}\) to the completed
regular sample and then transferring the resulting blue-witness predicate
back to the already fixed blue realization through the common injection,
recorded table, and blue terminal values.  The transfer is justified by
\(\noderef{FixedSetOppositeColorTerminalCoverage}\), which is part of the
adaptive helper's support: every blue coordinate needed to evaluate the
opposite-color neighborhoods is either recorded in \(R_h\) or lies in the
full terminal set \(U_h\).  The blue present set in this order is only
constrained to be a subset of the projected blue pairs, of which there are at
most \(k^2\).  When \(u_B\ge u_R\), reveal the red terminal candidates first
and use the symmetric fixed set \(E_B^\ast\) for the later blue coordinates.

Fix \(u_R,u_B\) with \(u_R\ge u_B\).  The product law on the fixed terminal
set \(U_h\), used sequentially along the prefix-safe order, bounds the
probability inside \(C_h\) that \(\mathcal B_I\cap\mathcal R\) occurs with
these cardinalities by
\[
\binom{3\varepsilon_1k^2}{u_R}p^{u_R}
\binom{k^2}{u_B}p^{u_B}
(1-p)^{f(\ell_R,\ell_B)-10\varepsilon_1k^2-u_R-u_B}. \tag{4}
\]
Indeed, first choose the \(u_B\) present blue coordinates from at most \(k^2\)
projected blue pairs and require those blue Bernoulli values to be present.
After this first-color realization is fixed, (3) gives the fixed red
candidate set \(E_R^\ast\) of size at most \(3\varepsilon_1k^2\).  Choose the
\(u_R\) present red coordinates from \(E_R^\ast\) and require those red
Bernoulli values to be present.  Prefix safety from
\(\noderef{FixedSetExposureCellProductLaw}\) justifies this sequential use of
the product law: staged-openness of the coordinate currently being scanned is
decided before its own Bernoulli value is read.  Finally require every other
staged-open coordinate in the realized adaptive set \(T(\omega)\) to be
absent.  The exponent of \(1-p\) is at least the lower bound in (2) minus the
two present counts.  Coordinates in \(U_h\setminus T(\omega)\) and the
additional condition \(\mathcal R\) can only remove branches from this count,
so they do not increase the upper bound.

Factor the absence probability in (4):
\[
\binom{3\varepsilon_1k^2}{u_R}
\left(\frac{p}{1-p}\right)^{u_R}
\binom{k^2}{u_B}
\left(\frac{p}{1-p}\right)^{u_B}
(1-p)^{f(\ell_R,\ell_B)-10\varepsilon_1k^2}. \tag{5}
\]
Set
\[
a=f(\ell_R,\ell_B)-10\varepsilon_1k^2.
\]
The eventual-comparisons package, together with
\[
p=\noderef{TwoBiteEdgeProbability}\!\left(\frac12,n\right),
\]
gives \(0\le p\le1/2\), hence \(p/(1-p)\le2p\).  We now split according to
the sign of \(a\).  If \(a<0\), then the desired bound for the fixed
\((u_R,u_B)\) is immediate: every conditional probability is at most \(1\),
while
\[
8\sqrt{\varepsilon_1}\,pk^2-pa\ge0
\]
because \(p\ge0\), \(\varepsilon_1\ge0\), and \(a<0\).  Hence
\(\exp(8\sqrt{\varepsilon_1}\,pk^2-pa)\ge1\).

It remains to treat the case \(a\ge0\).  Since \(0\le p\le1/2\), the base
\(1-p\) lies in \((0,1]\).  The standard inequality
\((1-p)^a\le\exp(-pa)\) is applicable precisely because \(a\ge0\), and gives
\[
(1-p)^{f(\ell_R,\ell_B)-10\varepsilon_1k^2}
\le
\exp\!\left(-p\bigl(f(\ell_R,\ell_B)-10\varepsilon_1k^2\bigr)\right).
\tag{6}
\]
Using the standard binomial estimate
\(\binom Ax\le(eA/x)^x\), with value \(1\) when \(x=0\), the two binomial
factors in (5) are at most
\[
\left(\frac{6e\varepsilon_1pk^2}{u_R}\right)^{u_R},
\qquad
\left(\frac{2epk^2}{u_B}\right)^{u_B}.
\tag{7}
\]
Since \(0\le\varepsilon_1\le1\) and \(u_B\le u_R\), replacing the first
\(\varepsilon_1\) by \(\sqrt{\varepsilon_1}\) and inserting
\(\sqrt{\varepsilon_1}\) into the second factor only increases the product.
The continuous maximum of \((A/x)^x\) is attained at \(x=A/e\), so (7) is
bounded by
\[
\exp(8\sqrt{\varepsilon_1}\,pk^2). \tag{8}
\]
In the nonnegative case \(a\ge0\), combining (5), (6), and (8) gives
\[
\Pr(\mathcal B_I\cap\mathcal R\cap\{|P_R|=u_R,|P_B|=u_B\}\mid C_h)
\le
\exp\!\left(
8\sqrt{\varepsilon_1}\,pk^2
-p\bigl(f(\ell_R,\ell_B)-10\varepsilon_1k^2\bigr)
\right). \tag{9}
\]
The negative case was handled above by the trivial probability bound, so
(9) holds for every sign of \(a\).
The case \(u_B\ge u_R\) is identical after swapping the colors, so (9) holds
for every ordered pair \((u_R,u_B)\).

Summing over the at most \(k^4\) ordered pairs gives, uniformly in \(h\),
\[
\Pr(\mathcal B_I\cap\mathcal R\mid C_h)
\le
k^4
\exp\!\left(
8\sqrt{\varepsilon_1}\,pk^2
-p\bigl(f(\ell_R,\ell_B)-10\varepsilon_1k^2\bigr)
\right).
\]
Since the large-\(n\) hierarchy gives \(k>0\), \(k^4=\exp(4\log k)\).  The
history-uniform averaging clause of
\(\noderef{FixedSetExposureCellProductLaw}\) transfers this same bound from
the histories \(C_h\) to the projection-size cell \(S\).  This is exactly the
asserted conditional exposure-cell bound.
\end{proof}
